% =====================================================================
%  CHƯƠNG 7 — TRIỂN KHAI THIẾT BỊ BIÊN  | Tuần 7-8 | Phụ trách: Nhật (tối ưu) + Đức (deploy)
% =====================================================================
\chapter{Triển Khai trên Thiết Bị Biên}
\label{ch:trienkhai}

\section{Nền Tảng Phần Cứng}
\wip{Raspberry Pi 5 (ARM Cortex-A76, RAM); giao tiếp camera; lý do chọn \autocite{ultralytics2026rpiguide}.}

\section{Tối Ưu Mô Hình}
\subsection{Xuất ONNX}
\wip{Chuyển mô hình sang ONNX; ONNX Runtime trên ARM.}
\subsection{Lượng Tử Hóa (Quantization)}
\wip{INT8 quantization; đánh đổi độ chính xác trước/sau nén \autocite{sonnara2025lightedge}.}
\subsection{Khả Chuyển Đa Nền Tảng}
\wip{Thiết kế độc lập nền tảng; nêu đường tối ưu thay thế (TensorRT trên Jetson) tuy phần
minh họa chính chạy trên Raspberry Pi 5 — bám đối tượng NC ở Chương~\ref{ch:tongquan}.}

\section{Đóng Gói và Triển Khai}
\wip{Service hóa pipeline; khởi động cùng hệ thống; (tùy chọn) container hóa.}

\section{Chạy Thử trên Raspberry Pi 5}
\wip{Cấu hình chạy; xác nhận hệ thống hoạt động end-to-end trên thiết bị biên; ảnh minh họa.}
